\section{Trabajo relacionado}
\label{sec:Related_work}


\paragraph{Multimodal Large Language Model.} Extender las capacidades multimodales desde imágenes estáticas hacia secuencias de video dinámicas introduce complejidad adicional, lo que exige que los modelos posean mayores capacidades para modelar dependencias de largo alcance y comprender eventos \cite{li2023videochat,liu2024kangaroo,zohar2024apollo,liu2024nvila,zhang2025videollama3,internvideo2}. Los avances recientes en MLLMs para video han introducido una variedad de estrategias innovadoras para enfrentar los desafíos de procesar y razonar eficientemente sobre entradas de video largas \cite{lin2023videollava,liu2024oryx,li2025llavast,shu2024videoxl,li2024videochatflash}. LongVILA \cite{chen2024longvila} propone un sistema de Multimodal Sequence Parallelism para modelado de contexto largo, habilitando entrenamiento e inferencia paralelos eficientes sobre contenido de video extendido. Sin embargo, la mayor parte de la investigación actual en video understanding sigue centrada en escenarios offline \cite{wang2024longllava,zeng2024timesuite,yan2024tpo,li2025videochatr1}, donde el modelo tiene acceso completo a la secuencia de video antes de la inferencia. Aunque esta configuración facilita el modelado semántico global, resulta insuficiente en escenarios de streaming, donde se requiere comprensión en tiempo real de escenas en evolución continua. Por lo tanto, el desarrollo de modelos diseñados específicamente para online video understanding es de importancia crítica.



\paragraph{Streaming Video Understanding.} En aplicaciones del mundo real, los usuarios esperan cada vez más que los MLLMs soporten procesamiento online e interacción en tiempo real. Esta demanda ha impulsado un interés creciente en la tarea de streaming video understanding. Recientemente, varios trabajos han explorado esta área emergente \cite{chen2024videollmonline,zhang2024flashvstream,yang2025svbench,di2025rekv,xiong2025streamchat,huang2024videochatonline}. Sin embargo, la mayoría de los enfoques existentes para streaming video understanding están diseñados principalmente para streaming dense video captioning \cite{chen2024videollmonline,wang2024mmduet,li2025lionfs,qian2025dispider,ding2025streammind}, centrándose únicamente en resumir contenido semántico a partir de frames visuales. Como resultado, tienen dificultades para manejar tareas esenciales como memory recall y percepción en tiempo real, que son críticas para un streaming video understanding integral. Además, en busca de eficiencia computacional, muchos métodos aplican compresión agresiva a secuencias de frames de video \cite{qian2024videostreaming, zhang2024internlmxco, yao2025timechatonline}, lo que los vuelve inadecuados para tareas complejas y dinámicas que requieren comprensión espaciotemporal fina y en tiempo real, como la conducción autónoma. Para abordar estas limitaciones, nuestro objetivo es desarrollar un enfoque más generalizable y práctico para online video understanding. Este enfatiza características espaciotemporales finas en el momento de la consulta y soporta almacenamiento persistente de memoria basado en eventos.
